\chapter{Arsitektur dan Teknologi Sistem}

\section{Arsitektur Sistem}
Sistem Koordinasi Kebencanaan BNPB dibangun menggunakan arsitektur \textit{event-driven microservices} yang terisolasi dalam batas container mandiri. Arsitektur sistem terdiri atas tiga tingkatan utama:
\begin{enumerate}
    \item \textbf{Lapisan Hulu (Upstream Mock Services):} Terdiri atas \texttt{mock-bmkg} (port 8081) dan \texttt{mock-pvmbg} (port 8082) yang mensimulasikan penyedia data instansi eksternal secara otonom dengan karakteristik latensi dan kegagalan yang dapat dikendalikan.
    \item \textbf{Lapisan Inti Koordinasi BNPB (Core Aggregation \& Persistence):} Terdiri atas \texttt{aggregator} (port 8083 internal) yang memegang kepemilikan eksklusif atas basis data kanonis \texttt{canonical-db} (PostgreSQL 16, port 5432) serta mempublikasikan event kanonis ke broker Redpanda (port 9092 / 19092).
    \item \textbf{Lapisan Hilir Klien \& Konsumen (Downstream Client API \& Event Consumers):} Terdiri atas \texttt{auth} (port 8084), \texttt{client-api} (port 8080), serta dua \textit{consumer} independen \texttt{dashboard-consumer} dan \texttt{notifier-consumer}.
\end{enumerate}

\begin{figure}[H]
\centering
\begin{tabular}{|c|}
\hline
\textbf{Diagram Topologi Arsitektur Terdistribusi BNPB} \\
\hline
\begin{minipage}{0.9\textwidth}
\vspace{0.3cm}
\centering
\texttt{[BMKG Mock :8081]} \quad $\xleftarrow{\text{Polling (X-BMKG-Key)}}$ \quad \textbf{[Aggregator :8083]} \quad $\xrightarrow{\text{Polling (Bearer)}}$ \quad \texttt{[PVMBG Mock :8082]} \\
\vspace{0.2cm}
$\downarrow$ \quad \textit{(Atomic Outbox TX)} \quad $\downarrow$ \\
\textbf{[PostgreSQL 16 (Canonical Store)]} \quad $\xleftarrow{\text{pgx Pool}}$ \quad \textbf{[Outbox Relay Worker]} \\
\vspace{0.2cm}
$\downarrow$ \quad \textit{(Publish HazardEvent)} \\
\textbf{[Redpanda Message Broker :9092 / topic: bnpb.hazard-events]} \\
\vspace{0.2cm}
$\swarrow$ \hspace{4cm} $\searrow$ \\
\texttt{[Dashboard Consumer]} \hspace{2.5cm} \texttt{[Notifier Consumer]} \\
\textit{(Group: bnpb-dashboard-group)} \hspace{1cm} \textit{(Group: bnpb-notifier-group)} \\
\vspace{0.3cm}
\hrule
\vspace{0.3cm}
\textbf{[Client API :8080]} \quad $\xleftarrow{\text{Internal REST}}$ \quad \textbf{[Aggregator]} \quad $\xleftarrow{\text{Verify JWT / Scope}}$ \quad \textbf{[Auth Service :8084]} \\
$\uparrow$ \textit{(HTTP GET /api/v1/hazards)} \\
\textbf{[Downstream Clients: Media (Summary) | Tim Lapangan (All/TTL 60s) | Internal]} \\
\vspace{0.3cm}
\end{minipage} \\
\hline
\end{tabular}
\caption{Diagram Arsitektur Sistem Terdistribusi BNPB Milestone 1}
\label{fig:arsitektur}
\end{figure}

\section{Tanggung Jawab Setiap Service}
Setiap service memiliki batas tanggung jawab (\textit{bounded context}) yang didefinisikan secara tegas tanpa adanya tumpang tindih logika:
\begin{itemize}
    \item \textbf{\texttt{mock-bmkg}:} Mensimulasikan data gempa bumi (\texttt{SeismicEvent}) dan peringatan tsunami (\texttt{TsunamiWarning}). Menegakkan autentikasi via header \texttt{X-BMKG-Key} dan merespons dengan latensi rendah (50--150 ms).
    \item \textbf{\texttt{mock-pvmbg}:} Mensimulasikan laporan gunung api (\texttt{VolcanicReport}). Menegakkan autentikasi via Bearer token, mendukung injeksi latensi tinggi (500 ms--3 s), simulasi kegagalan instan (\texttt{POST /admin/outage}), serta perubahan skema aditif saat sistem berjalan (\texttt{POST /admin/schema-version}).
    \item \textbf{\texttt{aggregator}:} Melakukan \textit{ingestion polling} berkala ke kedua sumber hulu secara paralel, melakukan normalisasi ke format kanonis \texttt{HazardEvent}, memperkaya koordinat gunung api via katalog statis, menyimpan data secara transaksional ke PostgreSQL, mendeteksi gangguan hulu dan menandai status degradasi, serta mengeksekusi \textit{Transactional Outbox Relay}.
    \item \textbf{\texttt{auth}:} Mengelola direktori klien, menerbitkan token akses JWT berumur 60 detik dengan klaim \textit{scope} terisolasi, serta mengelola rotasi \textit{refresh token} dengan deteksi \textit{replay revocation}.
    \item \textbf{\texttt{client-api}:} Menyediakan gerbang HTTP publik bagi pihak hilir. Bertanggung jawab memverifikasi tanda tangan JWT dan cakupan hak akses, memfilter atribut mentah untuk pengguna Media, serta meneruskan permintaan ke Aggregator dengan propagasi \texttt{X-Correlation-ID}.
    \item \textbf{\texttt{dashboard-consumer}:} Berlangganan topik event kanonis dengan grup \texttt{bnpb-dashboard-group}, melakukan deduplikasi berbasis \texttt{hazard\_id}, dan memperbarui metrik rekapitulasi bencana secara \textit{real-time}.
    \item \textbf{\texttt{notifier-consumer}:} Berlangganan topik event kanonis dengan grup \texttt{bnpb-notifier-group}, mengevaluasi ambang batas darurat kebencanaan (gempa $\ge 6.0$, tsunami SIAGA/AWAS, vulkanik AWAS), dan menerbitkan sinyal peringatan darurat.
\end{itemize}

\section{Protokol Komunikasi Antarservice}
Sistem menerapkan pemisahan yang ketat antara komunikasi sinkron dan asinkron:
\begin{itemize}
    \item \textbf{Komunikasi Sinkron (HTTP/REST):} Digunakan untuk operasi permintaan-respons (\textit{request-response}) berorientasi pembacaan data historis. Polling dari Aggregator ke mock upstream serta Client API ke Aggregator menggunakan HTTP/1.1 dengan payload JSON dan header \texttt{X-Correlation-ID}.
    \item \textbf{Komunikasi Asinkron (Kafka Protocol via Redpanda):} Digunakan untuk aliran diseminasi kejadian bencana ke banyak konsumen secara non-pemblokir (\textit{non-blocking pub-sub}). Aggregator menerbitkan pesan terkanonisasi ke topik \texttt{bnpb.hazard-events} menggunakan \texttt{hazard\_id} sebagai \textit{message key} guna menjamin keterurutan partisi.
\end{itemize}

\section{Kepemilikan Penyimpanan dan Batas Container}
Sesuai prinsip arsitektur \textit{microservices} dan klausul wajib spesifikasi Tugas Besar:
\begin{itemize}
    \item \textbf{Single Ownership:} Basis data relasional \texttt{canonical-db} dimiliki secara eksklusif oleh \texttt{aggregator}. Service lain (seperti \texttt{client-api} atau konsumen) secara arsitektural dilarang terhubung langsung ke port 5432 database.
    \item \textbf{Isolasi Jaringan Docker:} Port basis data tidak diekspos ke jaringan publik host. Komunikasi antara \texttt{aggregator} dan \texttt{canonical-db} berlangsung di dalam jaringan privat \texttt{bnpb-internal}.
    \item \textbf{Zero Shared Libraries:} Seluruh service Go dikembangkan secara mandiri tanpa menggunakan pustaka \textit{shared business logic} (\texttt{pkg/} atau \texttt{common/}) bersama. Masing-masing service memiliki berkas \texttt{go.mod} mandiri yang diorkestrasi via \texttt{go.work}.
\end{itemize}

\section{Alur Permintaan, Data, dan Event}
Alur pemrosesan data ujung-ke-ujung dirancang sebagai berikut:
\begin{enumerate}
    \item \textbf{Pengambilan (Ingestion):} Dua goroutine poller mandiri dalam Aggregator secara berkala memanggil \texttt{mock-bmkg} dan \texttt{mock-pvmbg} setiap 2 detik dengan jendela tumpang tindih (\textit{overlap}) 10 detik.
    \item \textbf{Normalisasi \& Transaksi:} Data dipetakan menjadi \texttt{HazardEvent}. Dalam satu transaksi ACID (\texttt{SaveHazardWithOutbox}), Aggregator menyimpan entitas bahaya ke tabel \texttt{hazards} sekaligus mencatat baris baru pada tabel \texttt{outbox\_events}.
    \item \textbf{Relay Outbox:} Goroutine latar belakang membaca baris outbox yang belum terkirim, mempublikasikannya ke Redpanda dengan header korelasi, lalu menandai baris outbox sebagai terkirim (\textit{marked as sent}).
    \item \textbf{Konsumsi \& Deduplikasi:} Konsumen \texttt{dashboard} dan \texttt{notifier} menarik pesan secara paralel, memeriksa tabel deduplikasi memori, memproses logika bisnis masing-masing, dan melakukan \textit{manual offset commit}.
\end{enumerate}

\section{Teknologi Terpilih dan Alasan Pemilihan}
Tabel~\ref{tab:teknologi} merangkum tumpukan teknologi yang dipilih beserta justifikasi teknisnya:

\begin{table}[htbp]
\centering
\small
\begin{spacing}{1.1}
\caption{Rangkuman Pemilihan Teknologi Sistem Terdistribusi BNPB}
\label{tab:teknologi}
\begin{tabularx}{\textwidth}{l p{4.2cm} X}
\toprule
\textbf{Kategori} & \textbf{Teknologi Terpilih} & \textbf{Alasan Pemilihan \& Kesesuaian Kebutuhan} \\
\midrule
Back-End & Go 1.22 & Dukungan konkurensi bawaan (\textit{goroutines} dan \textit{channels}) yang sangat efisien, konsumsi memori rendah, kompilasi cepat ke biner statis, dan menyeragamkan pengembangan seluruh service. \\
\addlinespace
Database & PostgreSQL 16 dengan JSONB & Menjamin transaksi ACID penuh untuk pola Transactional Outbox, sekaligus menyediakan fleksibilitas kolom JSONB dengan pengindeksan GIN untuk menyimpan atribut dinamis tanpa migrasi DDL. \\
\addlinespace
Komunikasi & HTTP/REST \& JSON & Standar industri yang transparan, mudah di-debug, mendukung interoperabilitas universal, dan menyederhanakan penegakan kontrak OpenAPI antarservice. \\
\addlinespace
Message Broker & Redpanda 24.1 (API Kafka) & Kompatibel penuh dengan ekosistem Kafka tanpa ketergantungan JVM atau Zookeeper, latensi deterministik berbasis C++/Raft, serta konsumsi sumber daya ringan dalam container. \\
\addlinespace
Autentikasi & JWT HS256 \& Refresh Token & Memungkinkan verifikasi klaim hak akses (\textit{scope}) secara terdesentralisasi di sisi Client API tanpa panggilan jaringan berulang ke Auth Service, dengan rotasi berkala untuk keamanan. \\
\addlinespace
Infrastruktur & Docker \& Docker Compose & Menyediakan batas proses yang terisolasi sempurna bagi setiap komponen serta memungkinkan orkestrasi seluruh topologi multi-service melalui satu perintah (\texttt{make up}). \\
\addlinespace
Dokumentasi & OpenAPI 3.0 & Mendefinisikan kontrak antarmuka API yang ketat dan terstandarisasi untuk mencegah pelanggaran skema data antartim. \\
\addlinespace
Pengujian & k6, curl, dan jq & k6 memungkinkan pengujian beban konkurensi tinggi ($\ge 50$ VUs) dengan pelaporan terstruktur; curl dan jq mempermudah otomatisasi skrip demonstrasi. \\
\addlinespace
Orkestrasi & GNU Make \& Fish/Bash & Menyeragamkan eksekusi skrip demonstrasi per problem (\texttt{make demo-p1} hingga \texttt{demo-p5}) agar pengujian dapat direproduksi secara instan. \\
\bottomrule
\end{tabularx}
\end{spacing}
\end{table}
